\paragraph{Completion of static and volumetric comparisons.}
The earlier standard-metric evaluation selected only the five planar temporal inputs, and the primary-table generator explicitly left ST-MTM unavailable on the four additions. We now compute ST-MTM on the unchanged Nucleon slice, all 21 reconstructed Benzene sections, all 201 Tangaroa volumes and all 102 Cylinder volumes: 325 additional layouts and complete 32768-row scalar profiles. All four methods now have SNS and TW on every input. TD remains undefined for static inputs; the Benzene column reports separately named adjacent-section displacement over 170 matches.

\paragraph{Fixed extension protocol.}
The common settings are uniform distance weights, $r=0.95$, $\lambda=0.5$, $\alpha=1$, minimum spacing $0.01$ in native coordinates, and start index zero. Nucleon uses base mode; the others propagate along the saved frame/section order using unchanged overlap correspondences. The initial translation uses the first frame's mean centroid coordinate. One reflection is fixed from the earliest informative anchor/reference covariance, with relative tolerance $10^{-12}$; no framewise reflection or score-based parameter selection occurs. Total leaf extent $K$ equals the existing width scale times the first-frame support measure. Tangaroa uses its saved physical voxel volumes; other inputs use their saved sample counts. The width targets are proportional to support size. This distinction preserves units and is necessary for a meaningful volumetric extension.

Three-dimensional source-centroid Euclidean distances enter the same OLO and distance objectives as in two dimensions. Genuine multifurcations are refined deterministically by original child ID into a binary right chain, using equal-level virtual saddles and empty new arcs. Every original leaf support and pairwise LCA scalar is checked unchanged. The refinement restricts the admitted orders; these runs are an explicit shared-input implementation extension, not original-author presets or byte-identical output reproduction.

\paragraph{Numerical solution and independent validation.}
The original SLSQP projection encountered numerical failure on the volumetric inputs. The completion instead solves the identical convex distance objective and linear spacing constraints with OSQP 1.1.3, recording primal, dual-stationarity and complementarity residuals. The proportional-width objective is nonnegative; feasible exact-target widths therefore certify the global minimum zero. Centered target intervals are used where feasible, otherwise a linear feasibility query supplies their placement. Unscaled relative-width losses verify these zero-objective certificates, avoiding misleading small residuals from a badly scaled width QP.

Independent saved-array checks recover all 325 profiles' labelled pairwise barriers and extremum values, reproduce SNS/TW/TD from retained centroids and anchors, reconstruct the distance-objective gradient and dual stationarity, and recount volumetric motion events. All new profiles have zero measured merge-level error. The 57 existing tests also pass. Historical results remain unchanged; source hashes, skeletons, landmarks, full-profile bindings, solver certificates and complete-precision summaries are retained under \path{evidence/stmtm_completion20261026/}. Complete scalar arrays reside in the corresponding reproducible result directory.

\paragraph{Motion-threshold scope.}
Tangaroa uses the retained method-specific occupied-range display threshold. Cylinder retains its original fixed full-canvas display threshold for every method, including the new ST-MTM run. Both use the same original within-dataset source extrema, overlap tracks, lag two and 2\%-of-domain-diagonal clear-event rule. Thus the method comparisons reuse exactly 6755 Tangaroa and 3196 Cylinder clear events, while the two datasets' display thresholds remain explicitly distinct. Static sections receive neither a temporal ME nor an artificial zero TD.
